\documentclass[10pt,a4paper]{amsart}
\usepackage[margin=19mm]{geometry}
\usepackage{amsmath,amssymb,mathtools}
\usepackage[scr=rsfs,cal=euler]{mathalfa}
\usepackage{xfrac}
\usepackage[unicode,hidelinks,bookmarks=false]{hyperref}
\newcommand{\ie}{i.e.\ }
\newcommand{\cf}{cf.\ }
\newcommand{\resp}{respectively\ }
\newcommand{\Subsection}{\subsection}
\providecommand{\nobreakdash}{\nobreak\hbox{-}}
\setlength{\emergencystretch}{2em}
\multlinegap0pt
\usepackage{bm}
\usepackage{bbm}
\usepackage{dsfont}
\usepackage{xspace}
\usepackage{cancel}
\usepackage{mathtools}
\usepackage{amsthm}
\makeatletter
\@ifundefined{theorem}{\newtheorem{theorem}{Theorem}}{}
\@ifundefined{definition}{\newtheorem{definition}[theorem]{Definition}}{}
\@ifundefined{lemma}{\newtheorem{lemma}[theorem]{Lemma}}{}
\makeatother
\title[Boundedness of Hausdorff-type operators with two-variable ke]{Boundedness of Hausdorff-type operators with two-variable kernels on Lebesgue spaces}
\author{A. R. Mirotin}
\date{Source: arXiv:2506.14333v1 (CC BY 4.0)}
\begin{document}
\maketitle

\noindent\emph{Source and licence.} Boundedness of Hausdorff-type operators with two-variable kernels on Lebesgue spaces. Version arXiv:2506.14333v1 (CC BY 4.0), distributed under the Creative Commons Attribution 4.0 International licence, as recorded in the arXiv licence metadata for this exact version and shown on the abstract page https://arxiv.org/abs/2506.14333v1. Full rights notice: licenses/arxiv-2506.14333-CC-BY-4.0.txt.\par\smallskip

\noindent\textit{Editorial scope.} Counted unit: the Theorem whose source label is \texttt{th:bdd\_Lqp\_two} in the article (source0.tex lines 175--182, source label \texttt{th:bdd\_Lqp\_two}), delivered together with the article's own complete original proof of it (source0.tex lines 184--240, one \texttt{proof} environment of 57 lines). No mathematics is written, repaired or re-derived: statement and proof are the article's own text line for line. Required context retained uncounted: df:general1\_two (lines 89--96); egree (lines 129--139); nuA(u) (lines 132--134); lm:IfA(u) (lines 141--146); r (lines 153--159). Every definition, display and label of those lines is retained unchanged. Omitted from this package and not redistributed: 1--88, 97--128, 135--140, 147--152, 160--174, 183--183, 241--438. Nothing inside a retained excerpt is omitted or altered: every retained body is delivered exactly as the article's lines are written, so no omission receipt of any kind accompanies this package. Citations: the retained excerpts carry no citation command at all, so no bibliography entry of the article is transcribed and no reference is redistributed. Reference adaptation: none was needed and none was made; every reference inside the delivered unit resolves inside it, so no unresolved reference or citation is produced. Printed slips kept literal: no printed word, symbol, hypothesis or number of the counted statement or of its proof was altered. Layout and class: the article's own class is \texttt{amsart} with packages including lineno, times, amsfonts, amssymb, latexsym, xcolor, comment; the wrapper is the lane's pinned \texttt{amsart} editorial wrapper at 10pt on A4, which restates the theorem-like environments the retained text needs and adds the article's own macro definitions that the retained text needs (none), with extra wrapper packages (bm, bbm, dsfont, xspace, cancel, mathtools). The delivered numbering is the unit's own. Assets: no retained span contains a figure environment, a graphics-inclusion command, a PDF-page inclusion, a table or a TikZ picture; no image file is copied into this package, the package declares no asset dependency and no image file is redistributed with it. Licence: version arXiv:2506.14333v1 of this article is distributed under the Creative Commons Attribution 4.0 International licence, as recorded in the arXiv licence metadata for this exact version and shown on the abstract page https://arxiv.org/abs/2506.14333v1. A full rights notice is delivered at licenses/arxiv-2506.14333-CC-BY-4.0.txt. Characters: every retained excerpt is pure ASCII, so the delivered engine, XeLaTeX with its default Latin Modern text font, reports no missing character.
\begin{definition}\label{df:general1_two}  Let  $S$ and  $S'$ be two sets, $(\Omega,\mu)$ denotes some measure space, and $A(u): S\to S'$ ($u\in  \Omega $) be some family of mappings\footnote{$A(u)$ do not assumed to be invertible.}.  Let  $\Phi(u,x)$ be a given  function on $\Omega\times S$ which is $\mu$-measurable for every fixed $x\in S$, and $V$ some Banach space.
 {\it A Hausdorff-type operator  with a two-variable kernel} acts on a  functions $f: S' \to V$  by the rule
\begin{equation}\label{general_two}
(\mathcal{H}_{\Phi,A,\mu}f)(x) =\int_{\Omega} \Phi(u,x)f(A(u)(x))d\mu(u), \ x\in S,
\end{equation}
%\end{comment}
provided the integral  converges in a suitable  sense.
\end{definition}

\begin{definition}\label{egree} %Let $(S, \mathcal{B}, \nu)$ be a   measure space.
 We say that a family $A(u): S\to S'$ ($u\in \Omega$) 
\textit{agrees with  measures $\nu$ and  $\nu'$ in a weak form}   if for each $\nu'$-measurable set $E\subset S'$ of finite measure and for every $u\in \Omega$ the set $A(u)^{-1}(E)$ is  $\nu$-measurable and
\begin{equation}\label{nuA(u)}
\nu(A(u)^{-1}(E))\le m(u)^{-1}\nu'(E)
\end{equation}
for some positive $\mu$-measurable function $m$ on $\Omega$.\footnote{In fact, $m(u)$  depends on $A(u)$.}

In the case where $(S,\nu)=(S',\nu')$ we say that $(A(u))_{u\in \Omega}$
\textit{agrees with   $\nu$  in a weak form}.
\end{definition}

\begin{lemma}\label{lm:IfA(u)}
Under the assumptions of  Definition \ref{egree} one has  for a function  $g\in L^1(\nu')$, $g(x')\ge 0$  that
\begin{equation}\label{IfA(u)}
\int_{S}g(A(u)(x))d\nu(x)\le m(u)^{-1}\int_{S'}g(x')d\nu'(x)
\end{equation}
\end{lemma}

\begin{equation}\label{r}
r=r(p,q):=
\begin{cases}
\frac{q}{q-p}, \mbox{ if } p\ne q,\\
\infty, \mbox{ if } p=q,
\end{cases} \ r(p,\infty):=1,
\end{equation}

\begin{theorem}\label{th:bdd_Lqp_two} Let $\infty>q\ge p\ge 1$,  or  $\infty=q> p\ge 1$, or $q=p= \infty$, the measure $\nu$ be   sigma-finite, and the family $(A(u))_{u\in \Omega}$ agrees with $\nu$ and $\nu'$  in a weak form.
If
\begin{equation}\label{Phi_qp}
 \|\Phi\|_{\mu,\nu,m}^{(p,q)}<\infty
\end{equation}
then the operator \eqref{df:general1_two}
 is bounded as an operator between $L^q(\nu')$ and  $L^p(\nu)$, and its norm does not exceed $\|\Phi\|_{\mu,\nu,m}^{(p,q)}$.
\end{theorem}

\begin{proof}
Let $p\ne \infty$. Using Minkowskii integral
inequality  we have for $1\le p<\infty$ and for an arbitrary $f\in L^q(\nu')$
\begin{eqnarray}\label{2_two}
\|\mathcal{H}_{\Phi, A,\mu}f\|_{L^p(\nu)}&=&\left(\int_{S} \left|\int_\Omega \Phi(u,x)f(A(u)(x))d\mu(u) \right|^pd\nu(x)\right)^{\frac 1p}\\ \nonumber
&\leq&
\int_\Omega\left(\int_{S}|\Phi(u,x)|^p|f(A(u)(x))|^pd\nu(x)\right)^{\frac 1p}d\mu(u).
\end{eqnarray}

First we assume that  $\infty>q\ge p\ge 1$.

Then by the Holder inequality and Lemma \ref{lm:IfA(u)}
\begin{align}\label{3_two}
\int_{S}|\Phi(u,x)|^p|f(A(u)(x))|^pd\nu(x)\\  \nonumber
&\hspace{-28mm}\le\left(\int_{S}|\Phi(u,x)|^{pr}d\nu(x)\right)^{\frac 1r} \left(\int_{S}|f(A(u)(x))|^{pr'}d\nu(x)\right)^{\frac 1{r'}}\\  \nonumber
&\hspace{-28mm}\le m(u)^{-\frac 1{r'}}\left(\int_{S}|\Phi(u,x)|^{pr}d\nu(x)\right)^{\frac 1r}
\left(\int_{S}|fx')|^{pr'}d\nu'(x')\right)^{\frac 1{r'}}
\end{align}
where $r$ is given by formula \eqref{r} and, as usual, $\frac 1r+\frac 1{r'}=1$.  Substituting this into \eqref{2_two} and taking into account that $pr'=q$ we have
\begin{eqnarray*}
\|\mathcal{H}_{\Phi, A,\mu}f\|_{L^p(\nu)}
\leq  \|\Phi\|_{\mu,\nu,m}^{(p,q)}\|f\|_{L^q(\nu')}
\end{eqnarray*}
and thus
\begin{eqnarray}\label{4_two}
\|\mathcal{H}_{\Phi, A,\mu}\|_{L^q(\nu')\to L^p(\nu)}
\leq  \|\Phi\|_{\mu,\nu,m}^{(p,q)}.
\end{eqnarray}

If $\infty=q> p\ge 1$, then in the previous argument the inequality \eqref{3_two} should be replaced by
\begin{align}\label{10_two}
\int_{S}|\Phi(u,x)|^p|f(A(u)(x))|^pd\nu(x)
\le\int_{S}|\Phi(u,x)|^{p}d\nu(x)\|f\|_{L^\infty(\nu')}^p
\end{align}
which is valid for every $u\in \Omega$ due to \eqref{nuA(u)}. Indeed, let 
$$
N':=\{x'\in S': |f(x')|>\|f\|_{L^\infty(\nu')}\},
$$
and
$$
N_u:=\{x\in S: |f(A(u)(x))|>\|f\|_{L^\infty(\nu')}\}.
$$
Then $\nu'(N')=0$, $N_u\subseteq A(u)^{-1}(N')$ and  therefore $\nu(N_u)=0$ for every $u\in \Omega$ by \eqref{nuA(u)}.
Thus, for every $u\in \Omega$ 
\begin{align}\label{20_two}
 |f(A(u)(x))|\le \|f\|_{L^\infty(\nu')}^p
 \end{align}
  for $\nu$-a.~e. $x\in S$
and therefore \eqref{10_two} holds for all $u\in \Omega$.

In view of \eqref{20_two} the proof in the case $q=p= \infty$ is almost obvious:
\begin{eqnarray*}
\|\mathcal{H}_{\Phi, A,\mu}f\|_{L^\infty(\nu)}&\leq& {\rm esssup}_{x\in S}\int_\Omega |\Phi(u,x)||f(A(u)(x))|d\mu(u)\\
&\leq & {\rm esssup}_{x\in S}\int_\Omega |\Phi(u,x)|d\mu(u)\|f\|_{L^\infty(\nu')}\\
&=&\|\Phi\|_{\mu,\nu,m}^{(\infty,\infty)}\|f\|_{L^\infty(\nu')}.
\end{eqnarray*}
\end{proof}

\end{document}
